%%%%%%%%%%%%%%%%%%%%%%%%%%%%%%%%%%%%%%%%%%%%%%%%%%%%%%%%%%%%%%%%%%%%%%%%%%%%%%%%
%   Copyright 2009, Oracle and/or its affiliates.
%   All rights reserved.
%
%
%   Use is subject to license terms.
%
%   This distribution may include materials developed by third parties.
%
%%%%%%%%%%%%%%%%%%%%%%%%%%%%%%%%%%%%%%%%%%%%%%%%%%%%%%%%%%%%%%%%%%%%%%%%%%%%%%%%

\section{Literals}
\seclabel{literals}

\note{We need to check whether the various types and values defined in this section are still correct.}

\begin{Grammar}
\emph{LiteralExpr} &::=& \texttt{(} \texttt{)}\\
&$|$& \emph{NumericLiteralExpr}\\
&$|$& \emph{CharLiteralExpr}\\
&$|$& \emph{StringLiteralExpr}\\
\end{Grammar}

Fortress provides boolean literals, \TYP{()} literal, character literals,
string literals, and numeric literals.  Literals are values; they do not
require evaluation.

\note{From here, copied from F1.0$\beta$.}

The literal \VAR{false} has type
\EXP{\TYP{BooleanLiteral}\llbracket\VAR{false}\rrbracket}.
The literal \VAR{true} has type
\EXP{\TYP{BooleanLiteral}\llbracket\VAR{true}\rrbracket}.
\note{And these are the only values of those types.}

The literal \EXP{()} is the only value with type \TYP{()}.
Whether any given occurrence of \TYP{()} refers to the value \EXP{()} or
to the type \TYP{()} is determined by context.

A character literal has type \TYP{Character}.
Each character literal consists of an abstract character in
\unicode~\cite{Unicode},
enclosed in single quotation marks (for example, `\txt{a}', `\txt{A}',
`\txt{\$}', `\txt{$\alpha$}', `\txt{$\oplus$}').
\marginnote{This is not true: Character literals include ``\b'' and
``LATIN CAPITAL A'', which do not ``consist of'' an abstract character enclosed in
single quotation marks.  I think it's better to remove most of this discussion and
direct the reader to the relevant section of Lexical Structure.
This section should be about {\bf semantics} of literals:
their values and types.}
The single quotes may be either true typographical ``curly'' single quotation marks
or a pair of ordinary apostrophe characters (for example, \txt{'a'}, \txt{'A'}, \txt{'\$'}, \txt{'$\alpha$'}, \txt{'$\oplus$'}).
See \secref{preprocessing-character-literals} for a description of how names of
characters may be used rather than actual characters within character literals,
for example `\txt{APOSTROPHE}' and `\txt{GREEK\_CAPITAL\_LETTER\_GAMMA}'.

A string literal has type \TYP{String}.
Each string literal is a sequence of \unicode\ characters enclosed in double
quotation marks (for example,
``\txt{Hello, world!}'' or ``\txt{$\pi\thinspace r^2$}'').
\marginnote{Again, I think this paragraph should mostly be eliminated and replaced with a pointer to the relevant section of Lexical Structure.

Here's a sentence I wrote as a replacement:
``A string literal (i.e., a sequence of characters enclosed in double quotation marks---see Section 5.10 for details) has type String.''}
The double quotes may be either true typographical ``curly'' double  quotation marks
or a pair of ``neutral'' double-quote characters (for example,
\txt{"Hello, world!"} or \txt{"$\pi\thinspace r^2$"}).
\secref{preprocessing-string-literals} also describes how names of
characters may be used rather than actual characters within string literals.
One may also use the escape sequences \txt{{\char'134}b} and \txt{{\char'134}t}
and \txt{{\char'134}n} and \txt{{\char'134}f} and \txt{{\char'134}r} as described in \cite{JLS}.
Two string literals of the same text denote one object: every evaluation of a
string literal with a given text, wherever it is written, yields the same
string, so that
\EXP{\hbox{\rm``\STR{ab}''} \sequiv \hbox{\rm``\STR{ab}''}} is \VAR{true}
(\secref{equivalence}).
A string that a program builds as it runs is another object than a literal of
the same text, unless the operation that builds it answers an existing string;
strings remain reference objects, whose equality of text is their \EXP{=}
(\chapref{objects}).
A concatenation of string literals is itself a string literal
(\secref{static-expr}).
\revision{revival-string-literals}{The Working Draft of February 2011 did not
say whether two string literals of the same text are one object.
The revival decided on 8~November 2026 that they are, on both implementations,
and that a concatenation of two string literals is folded into one literal.}


\note{We may want to eliminate the distinction between simple and compound numerals,
and just distinguish them by whether they have a radix point or not, when necessary
(i.e., ``Numerals may have a radix point.'').  Again,
that distinction is primarily a syntactic one, and can be discussed there if appropriate.}
Numeric literals in Fortress are referred to as \emph{numerals},
corresponding to various expressible numbers.
Numerals may be either \emph{simple} or \emph{compound}
(as described in \secref{numerals}).

Every numeral has the type \TYP{Numeral}, which extends \TYP{Number}, through
one of the two types that \TYP{Numeral} comprises:
a numeral containing only digits has type \TYP{NaturalNumeral},
and a numeral containing a radix point has type \TYP{RadixPointNumeral}.
The value of a numeral holds four numbers:
the number $n$ of its digits, leading zeros among them;
the number $m$ of its digits after the radix point, which is $0$ for a
\TYP{NaturalNumeral};
its radix $r$;
and $v$, the value of its digits, with the radix point deleted, interpreted in
radix $r$.
\Library\ declare these numbers as the getters \VAR{digits},
\VAR{fractionDigits}, \VAR{radix} and \VAR{digitValue} of \TYP{Numeral}, the
first three of type \EXP{\mathbb{Z}32} and the last of type \EXP{\mathbb{Z}}.
They are fields of the numeral's value, not static parameters of its type:
the three numeral types have no static parameters.
So \EXP{3.125} is a \TYP{RadixPointNumeral} whose four numbers are $4$, $3$,
$10$ and $3125$, and \EXP{007} a \TYP{NaturalNumeral} whose four numbers are
$3$, $0$, $10$ and $7$.
\note{Can the type \EXP{\TYP{NaturalNumeral}\llbracket{}n, 10, v\rrbracket}
have any value other than what you get from evaluating the corresponding literal?

What is the relationship between \EXP{\TYP{Literal}\llbracket{}v\rrbracket}
and \EXP{\TYP{NaturalNumeral}\llbracket{}n, 10, v\rrbracket}?

We need to describe the \TYP{Numeral} type hierarchy.
All the numeric/numeral types should be defined and their relations to each other given.
And a reference to where that is done should be appear in this section (probably at the beginning of the section).
}

\note{Ought we discuss digit-group separators here?
Or is that just part of lexical structure?
(I think we intend that a numeral consisting of digits and digit-group separators has
type \EXP{\TYP{NaturalNumeral}\llbracket{}n,10,v\rrbracket}
for appropriate values of \VAR{n} and \VAR{v}.)}

The radix of a numeral is ten, unless its digits are followed by an underscore
and a radix indicator (\secref{numerals}), which give its radix $r$; such a
numeral has the same type as one without an indicator.
So \EXP{\mathtt{deadc}_{16}} is a \TYP{NaturalNumeral} whose four numbers are
$5$, $0$, $16$ and $912092$.
\revision{revival-numerals}{The Working Draft of February 2011 gave a numeral
the types \EXP{\TYP{NaturalNumeral}\llbracket{}n,10,v\rrbracket},
\EXP{\TYP{NaturalNumeralWithExplicitRadix}\llbracket{}n,r,v\rrbracket},
\EXP{\TYP{RadixPointNumeral}\llbracket{}n,m,10,v\rrbracket},
\EXP{\TYP{RadixPointNumeralWithExplicitRadix}\llbracket{}n,m,r,v\rrbracket},
\EXP{\TYP{Literal}\llbracket{}v\rrbracket} for a numeral without leading zeros,
and \EXP{\TYP{Numeral}\llbracket{}n,m,r,v\rrbracket}, whose static arguments are
the four numbers.
The revival decided on 8~November 2026 that the four numbers are fields of the
numeral's value: the two types with an explicit radix fold into the radix, and
\TYP{Literal} into the digit count, which tells whether a numeral has leading
zeros.
So the team's draft notes above are answered: two numerals of one kind have one
type, which holds every such numeral, and \TYP{Numeral} comprises the two
numeral types.}


A numeral is converted where it stands, into the type that its place fixes,
and not before, because, as in common mathematical usage,
we expect numerals to be polymorphic.
For example, consider the numeral \EXP{3.1415926535897932384};
converting it to a floating-point number before its place is known
may lose precision.
If that numeral is used where a floating-point interval is expected,
it is better to convert it directly to an interval.
Therefore, numerals have their own types as described above.
This approach allows library designers to decide how numerals should
interact with other types of objects by defining coercion operations
(see \secref{coercion} for an explanation of coercion
in Fortress).
\Library\ define coercions
from numerals to integers (for simple numerals)
and rational numbers (for compound numerals).
In \library, each of the eight number types (\secref{basic-rationals}) declares
a coercion from \TYP{NaturalNumeral}, and \EXP{\mathbb{Q}}, \EXP{\mathbb{R}64}
and \EXP{\mathbb{R}32} declare one from \TYP{RadixPointNumeral}; the conversion
is exact into the five integer types (\secref{basic-integers}) and
\EXP{\mathbb{Q}}, and rounds once, to the nearest value of the type, into
\EXP{\mathbb{R}64} and \EXP{\mathbb{R}32}.

A place \emph{fixes a type} for a numeral when the type declared there declares
a coercion from the numeral's type, as each number type of \library\ does and a
type of a program's own may, or is a numeral type itself.
The declared type of a variable, a field or a parameter, and the declared
return type of a functional, are such places, the parameter types of the
declaration that a call invokes among them (\secref{resolving-coercion}).
There the numeral is converted into that type, from its exact value:
so \EXP{q\COLON\mathbb{Q} = 0.1} binds the rational \EXP{1/10},
\EXP{r\COLON\mathbb{R}64 = 0.1} the floating-point number nearest to it, and
\EXP{n\COLON\mathbb{N}32 = 5} the natural number 5.
It is a static error if the numeral's value lies outside the range of that
type, or rounds to an infinity there.
Where nothing fixes a type, as where a variable is declared without a type, or
a parameter has the type \TYP{Any}, \TYP{Object} or \TYP{Number} or a type
parameter declared without a bound, the numeral is converted into its
\emph{default}: a numeral without a radix point into \EXP{\mathbb{Z}32}, or into
\EXP{\mathbb{Z}64} or \EXP{\mathbb{Z}} when its value does not fit in the
narrower type, and a numeral with a radix point into \EXP{\mathbb{R}64}.
So \EXP{x = 3} binds \VAR{x} as a \EXP{\mathbb{Z}32},
\EXP{\VAR{bigN} = 3000000000} binds \VAR{bigN} as a \EXP{\mathbb{Z}64} and
\EXP{y = 0.5} binds \VAR{y} as an \EXP{\mathbb{R}64};
\EXP{x \sequiv z} is \VAR{true} for \EXP{z\COLON\mathbb{Z}32 = 3};
and for \EXP{\KWD{object}\;\TYP{Cell}\llbracket{}T\rrbracket(x\COLON{}T)},
\EXP{\TYP{Cell}(3)} is a \EXP{\TYP{Cell}\llbracket\mathbb{Z}32\rrbracket}.
\secref{numeral-default} says where else a numeral takes its default, as where
several numerals fix one type parameter, and how a numeral whose conversions
tie is read.
The default applies to a numeral, and to a static expression of numerals once
it is folded (\secref{static-expr}), never to an expression of a numeral type.

A type of a program's own receives a numeral through its own coercion from the
numeral's type, such as \EXP{\KWD{coerce}(x\COLON\TYP{RadixPointNumeral})},
whose body reads the four numbers.
A numeral at a place whose type is a type of a program's own that declares no
coercion from the numeral's type is a static error: coercion does not chain
(\secref{coercion}), so the numeral is not converted into its default first.
The coercion sees the four numbers of a numeral as written: a negative constant
reaches it as the prefix operator \EXP{-} applied to the coercion's result for
the natural numeral (\secref{static-expr}), so a type without that operator
refuses it.
At each place where such a numeral stands, the coercion is applied once, not at
each evaluation of the place, so that the numeral there denotes one object, as
a string literal does: two evaluations of one place yield the same object, and
two places yield two objects.
A numeral type is written in a program only as the type of a parameter, of a
coercion or of an operator, never as the type of a variable, a field or a
functional's result, nor as a static argument; so the constants \VAR{pi} and
\VAR{e} of \library\ are values of \EXP{\mathbb{R}64}.
Inside the body of a coercion, two values of a numeral type are equal, and
strictly equivalent, when their values $v/r^m$ are equal: there \EXP{2.5} and
\EXP{2.50} are strictly equivalent, and so are \EXP{0.1_{2}} and \EXP{0.5}.
\revision{revival-numerals}{The Working Draft of February 2011 said that
numerals are not directly converted to any of the number types, and so gave a
numeral where nothing fixes a type its own type.
The revival decided on 8~November 2026 that a numeral is converted where it
stands: where its place fixes a type, into that type, and elsewhere into the
default above; that a type of a program's own receives a numeral only through
its coercion from the numeral's type, once for each place; and that the numeral
types are written only as parameter types.
The text this replaces, the revival's own of 29~September 2026
(\secref{revival-numeral-type}), named the type of an integer numeral
\TYP{IntLiteral}, the name that the compiled path's own library gives it.}


\marginnote{What is an ``integer'' in this context?}
In Fortress, dividing two integers using the $/$ operator produces
a \emph{rational number};\marginnote{Not if the second integer is 0.}
this is true regardless of whether
the integers are of type $\mathbb{Z}$ (or \txt{ZZ}),
$\mathbb{Z}64$ (or \txt{ZZ64}), $\mathbb{N}32$ (or \txt{NN32}), or
whatever.  Addition, subtraction, multiplication, and division of rationals
are always exact; thus values such as \EXP{1/3} are represented exactly in
Fortress.

\note{We should give an example with a radix specifier.}
Numerals containing a radix point are actually rational
literals; thus \EXP{3.125} has the rational value \EXP{3125/1000}.
The quotient of two integer literals is a static expression
(described in \secref{static-expr}) whose value is rational.
Similarly, a sequence of digits with a radix point followed by the symbol
$\times$ and an integer literal raised to an integer literal,
such as \EXP{6.0221415 \times 10^{23}}
is a static expression whose value is rational.
If such static expressions are mentioned as part of a floating-point
computation, the compiler performs the rational arithmetic exactly and then
converts the result to a floating-point value, thus incurring
at most one floating-point rounding error.
But in general rational computations may also be performed at run time,
not just at compile time.
\revision{revival-numerals}{The revival decided on 8~November 2026 that an
expression of numerals and operators alone is folded exactly into a signed
constant before the static check, on both implementations
(\secref{static-expr}): so \EXP{x = 7/2} binds the \EXP{\mathbb{Q}} \EXP{7/2},
\EXP{1/2 = 0.5} is \VAR{true}, and the constant is rounded at most once where
a floating-point type is declared.}
\note{This specifies too much about the implementation:
I don't think we should mention a compiler at all, except perhaps as informal discussion.
In particular, wouldn't it be correct if this arithmetic were actually deferred and
performed exactly at run time?}


A rational number can be thought of as a pair of integers $p$ and $q$
that have been reduced to ``standard form in lowest terms'';
that is, $q > 0$ and there is no nonzero integer $k$
such that $p \over k$ and $q \over k$ are integers
and $\left| {p \over k} \right| + \left| {q \over k} \right|
	< \left| p \right| + \left| q \right|$.
The type $\mathbb{Q}$ includes all such rational numbers.
Relaxing the requirement $q>0$ to $q \geq 0$ adds three more values, which
$\mathbb{Q}$ also includes: \EXP{1/0} and \EXP{-1/0}, the infinite rationals,
and \EXP{0/0}, the indefinite rational.
\marginnote{Is there any subtyping or coercion relationship between \EXP{\mathbb{Q}} and
\EXP{\mathbb{Q}*}?
If \EXP{\mathbb{Q}} is a subtype of \EXP{\mathbb{Q}*}
then static checking would normally disallow the assignment,
so is this a relaxation to that rule?  If so, the static checker needs to know about it.}
The advantage of these values is that $\mathbb{Q}$ is closed under the rational
operations $+$, $-$, $\times$, and $/$.
In ASCII, \EXP{\mathbb{Q}} is written as \txt{QQ}.
See \secref{basic-rationals} for its operations.
\revision{revival-literals}{The Working Draft of February 2011 described
\EXP{\mathbb{Q}}, \EXP{\mathbb{Q}^*} and \EXP{\mathbb{Q}^{\#}} as three types,
said that a value of type \EXP{\mathbb{Q}^*} assigned to a variable of type
\EXP{\mathbb{Q}} throws a \TYP{DivisionByZero} at run time if it is \EXP{1/0}
or \EXP{-1/0}, gave the ASCII names \txt{QQ\_star} and \txt{QQ\_splat}, and
pointed to \secref{advanced-rationals} for the definitions.
The three are one type now, as in \library\ (\secref{basic-rationals}), so no
such assignment converts or throws; that is also the answer to the margin
note's question whether \EXP{\mathbb{Q}} and \EXP{\mathbb{Q}^*} are related by
subtyping or by coercion.}



\subsection{Pi}
\seclabel{pi}

The object named \EXP{\pi} (or \txt{pi}) may be used to represent the ratio
of the circumference of a circle to its diameter rather than a specific
floating-point value or interval value.
In Fortress, \EXP{\pi} has type
\EXP{\TYP{RationalValueTimesPi}\llbracket\VAR{false},1,1\rrbracket}.
When used in a floating-point computation, it becomes
a floating-point value of the appropriate precision; when used in an interval
computation, it becomes an interval of the appropriate precision.

\subsection{Infinity and Zero}

The object named \EXP{\infty} has type \EXP{\TYP{ExtendedIntegerValue}\llbracket\VAR{true},0,\VAR{true}\rrbracket}.
One can negate \EXP{\infty} to get a negative infinity.

\note{What about spaces?  That is, is ``- 0'' treated the same as ``-0'', for example?}
Negating the literal \EXP{0} produces a special negative-zero object,
which refuses to participate in static arithmetic
(discussed in \secref{static-expr}).  It
has type \TYP{NegativeZero}.  The main thing it is good for is coercion to
a floating-point number (discussed in
\chapref{conversions-coercions}).
(Negating any other zero-valued expression simply produces zero.)
